\section{Experimental setup}\label{sec:setup}

\subsection{Corpora and stream construction}
All data are public. Human-written traffic comes from the unified release \cite{ds_seven} of the seven corpora curated by Champa et al.\ \cite{champa2024}: Enron and Ling-Spam with SpamAssassin (1999--2004), TREC-05/06/07 (2005--2007) and CEAS-08 (2008), 203{,}017 messages in total. Their positive class is the union of phishing and spam, which we refer to as \emph{malicious} and discuss in Section~\ref{sec:validity}. LLM-written traffic comes from two independent releases: \emph{GPTGen} \cite{ds_gptgen}, 3{,}746 English phishing and legitimate messages generated with OpenAI models, and \emph{SpearGen} \cite{ds_spear}, 2{,}798 personalised spear-phishing messages generated by an OpenAI model from a name, a job title and a list of recent activities. We remove inline sender addresses from GPTGen (they occur in no other corpus and would act as an artefact), normalise whitespace, prefix human messages with their subject line, discard human messages of 40 characters or fewer and SpearGen outputs of 80 characters or fewer and remove exact duplicates across all sources. Disjoint pools are then drawn with a fixed seed (Table~\ref{tab:pools}); no text appears in more than one pool.

\begin{table}[t]
\centering
\caption{Disjoint pools (legitimate/malicious counts). Each realisation samples its segments from these pools.}
\label{tab:pools}
\scriptsize
\setlength{\tabcolsep}{2.5pt}
\begin{tabular}{@{}lrr>{\raggedright\arraybackslash}p{9cm}@{}}
\toprule
Pool & Legit. & Malic. & Sources (legit./malic.) \\
\midrule
warm-up & 2295 & 1705 & Assassin 433/175; Enron 1605/1488; Ling 257/42 \\
segment A & 4649 & 3351 & Assassin 851/344; Enron 3301/2909; Ling 497/98 \\
segment B & 4474 & 3526 & TREC-05 2119/1455; TREC-06 781/256; TREC-07 1574/1815 \\
segment C & 3582 & 4418 & CEAS-08 3582/4418 \\
segment D & 4400 & 3300 & CEAS-08 607/0; GPTGen 900/1500; SpearGen 0/1800; TREC-05 1400/0; TREC-06 507/0; TREC-07 986/0 \\
segment E (orig.) & 2700 & 0 & CEAS-08 471/0; TREC-05 1036/0; TREC-06 380/0; TREC-07 813/0 \\
segment E (T1) & 300 & 1500 & CEAS-08 0/119; GPTGen 300/500; SpearGen 0/500; TREC-05 0/149; TREC-06 0/34; TREC-07 0/198 \\
stress pool & 0 & 600 & CEAS-08 0/51; GPTGen 0/200; SpearGen 0/200; TREC-05 0/62; TREC-06 0/13; TREC-07 0/74 \\
\bottomrule
\end{tabular}
\end{table}


Each stream realisation (seed $r=0,\dots,9$) samples, without replacement, the five segments in a fixed order, shuffling messages inside each segment: A (6{,}000 human messages, 1999--2004), B (6{,}000, TREC), C (6{,}000, CEAS-08), D (6{,}000: LLM phishing from GPTGen and SpearGen at 30\% prevalence, mixed with GPTGen legitimate messages and held-out human legitimate traffic) and E (4{,}000: a paraphrase-evasion campaign of 1{,}500 T1-attacked messages, about 1{,}250 of them malicious and drawn from the two LLM corpora and from human sources, the rest a paraphrased legitimate control, mixed with 2{,}500 human legitimate messages). Malicious prevalence is 42\%, 44\%, 55\%, 30\% and 31\% in A to E. Rewrites that stayed below the similarity budget after three tries are rejected and sent unmodified (3.4\% of attempts, about 45 messages per E segment). The boundaries A$\to$B$\to$C are natural temporal and source shifts; C$\to$D is a generator shift; D$\to$E is an adversarial shift. The stream is processed in $b=250$-message batches with label delay $D=2$ batches. The warm-up window (4{,}000 messages from the 1999--2004 corpora) is the only data used to fine-tune the encoder and to fit the standardisation; every initial detector is trained on it alone. A held-out stress pool of phishing messages from all three malicious sources is never streamed and is used only in Section~\ref{sec:stress}.

\subsection{Leakage and artefact audit}
Generated corpora can encode their label in templates rather than semantics \cite{cetlenbik2026}. Before any experiment we therefore audited every source$\times$class cell for template collapse after entity masking (URLs, addresses, numbers and capitalised spans), near-duplicates (pairwise TF-IDF cosine on up to 1{,}500 sampled messages), length balance, a zero-learning opening-trigram lookup rule, and a human-versus-LLM authorship probe on legitimate mail. The results (Section~\ref{sec:audit}) shape the evaluation: SLIDE and all baselines are scored separately on LLM-written legitimate messages so that a detector that learns ``LLM style'' instead of ``malicious intent'' is exposed by its false-positive rate.

\subsection{Compared methods}
Table~\ref{tab:methods} lists the ten methods, grouped by the supervision they consume. Static detectors never see stream labels; the fully supervised adaptive detectors receive every label after the delay $D$, whereas label-on-alarm detectors buy labels only when their own drift test fires. To isolate the effect of the adaptation policy, every head-based method uses the same representation $\Phi$ and the same head architecture; the baselines are trained with plain class-weighted cross-entropy, whereas SLIDE uses Eq.~\eqref{eq:ltotal}.

\begin{table}[t]
\centering
\caption{Compared detectors. ``Labels'' is the supervision consumed during the stream.}
\label{tab:methods}
\scriptsize
\setlength{\tabcolsep}{3pt}
\begin{tabular}{@{}>{\raggedright\arraybackslash}p{2.8cm}>{\raggedright\arraybackslash}p{9.5cm}l@{}}
\toprule
Method & Description & Labels \\
\midrule
TFIDF-LR & word 1--2-gram TF-IDF + logistic regression, warm-up only & none \\
Static-FT & fine-tuned DistilRoBERTa classifier, frozen & none \\
Static-MLP & head $h_\theta$ on $\Phi$, cross-entropy, frozen & none \\
Online-TFIDF-SGD & hashed 1--2-grams, logistic SGD, \texttt{partial\_fit} per delayed batch & all \\
Online-MLP & $h_\theta$ fine-tuned for 3 epochs on each delayed batch & all \\
Periodic-Retrain & $h_\theta$ refit on the last 2{,}000 labels every 8 batches & all \\
ADWIN-Retrain & ADWIN ($\delta=0.002$) on the 0/1 error stream \cite{bifet2007}; refit on last 1{,}000 labels & all \\
PH-Retrain & Page--Hinkley (threshold 50) on the error stream; refit on last 1{,}000 labels & all \\
KSWIN-Score & KSWIN ($\alpha=0.005$) \cite{raab2020} on the score stream; on alarm buys $n_\ell$ labels, plain refit & on alarm \\
\textbf{SLIDE} & sliced two-gate test on $\Phi$; on alarm buys $n_\ell$ labels; $\mathcal{L}_{\mathrm{total}}$ with replay & on alarm \\
\bottomrule
\end{tabular}
\end{table}


Three controls separate the timing of label purchases from the update rule. KSWIN-Score+SLIDE-update keeps the score-based KSWIN trigger but updates with Eq.~\eqref{eq:ltotal}. The two budget controls apply the SLIDE update after exactly four purchases of $n_\ell$ labels, the same 2{,}000 labels that SLIDE buys on this stream: Budget-Random places the purchases at four batches drawn uniformly at random per realisation, and Budget-Oracle places each one at the first batch by which $n_\ell$ messages of a new segment have arrived. The oracle uses knowledge of the segment boundaries that no deployed detector has; it serves as an upper reference for alarm timing.

\subsection{Hyper-parameters and statistical protocol}
Hyper-parameters (Table~\ref{tab:hyper}) were fixed before the stream experiments from defaults in the cited works and from the warm-up window; none was selected on stream data. Sensitivity to $\epsilon$, $\alpha$, $n_\ell$ and $K$ is reported in Section~\ref{sec:sens}.

\begin{table}[t]
\centering
\caption{Hyper-parameters (fixed before the stream experiments).}
\label{tab:hyper}
\scriptsize
\begin{tabular}{@{}llr@{}}
\toprule
Component & Parameter & Value \\
\midrule
Encoder fine-tuning & epochs / batch / learning rate & 2 / 16 / $2\times10^{-5}$ \\
 & warm-up / schedule & 6\% linear / linear decay, AdamW, wd $0.01$ \\
 & max.\ tokens $L$ & 128 \\
Head $h_\theta$ & hidden width / dropout & 256 / 0.1 \\
 & epochs / batch / learning rate & 30 / 256 / $10^{-3}$ (AdamW) \\
Objective (Eq.~\ref{eq:ltotal}) & $\lambda_r$ / $\lambda_a$ / $\lambda_p$ & 1.0 / 0.5 / $10^{-3}$ \\
 & $\epsilon$ / PGD steps $P$ & 1.0 / 3 \\
Drift test & directions $K$ / level $\alpha$ / null draws $B$ & 64 / 0.01 / 100 \\
 & $n_R$ / $n_W$ & 1{,}000 / 500 \\
Policy & alarm window $n_\ell$ / memory per class $c$ & 500 / 1{,}000 \\
Stream & batch $b$ / delay $D$ / realisations & 250 / 2 / 10 \\
Attack T1 & accepted similarity / tries / prefix & $\ge 0.70$ / 3 / 600 chars \\
\bottomrule
\end{tabular}
\end{table}


We report precision, recall and F1 of the malicious class together with the Matthews correlation coefficient (MCC), the area under the precision--recall curve (PR-AUC, average precision), the false-positive rate (FPR), the Brier score and the number of stream labels consumed. Stream-level metrics pool all prequential predictions of a realisation; we report the mean over the ten realisations with a 95\% Student-$t$ confidence interval. Following Dem\v{s}ar \cite{demsar2006}, the primary test unit is the realisation: we apply the Friedman test across methods with realisations as blocks, the Nemenyi critical difference at $\alpha=0.05$, and two-sided Wilcoxon signed-rank tests of SLIDE against each competitor on the ten paired stream-level MCC values with Holm correction over the competitors. With ten pairs the smallest exact two-sided Wilcoxon $p$-value is $2/2^{10}\approx0.002$; after Holm correction over nine competitors the smallest attainable adjusted value is $9\times0.00195=0.0176$, which is therefore the strongest evidence this design can express. As a secondary, finer-grained analysis we apply the same test to the $10\times112$ paired batch-level MCC values; because consecutive batches are not independent we treat these $p$-values as descriptive. Effect sizes are matched-pairs rank-biserial correlations $r_{rb}\in[-1,1]$.

\subsection{Computing environment and reproducibility}\label{sec:repro}
All experiments ran on one Intel Xeon processor at 2.10\,GHz with four cores, 15\,GB of RAM and no GPU. The software stack is Python~3 with PyTorch 2.14 (CPU build), Hugging Face \texttt{transformers} 5.20, scikit-learn 1.9, SciPy 1.18 and \texttt{river} 0.26. The encoder is \texttt{distilroberta-base} \cite{sanh2019,liu2019roberta}, the surprisal model \texttt{distilgpt2}, the paraphraser a T5-base model fine-tuned for paraphrasing \cite{hf_paraphraser}, and the similarity encoder \texttt{all-MiniLM-L6-v2}. Encoder fine-tuning on the warm-up window took 25 minutes; the 240 stream runs of Sections~\ref{sec:main}--\ref{sec:sens} took 15--30\,s each once representations were cached.

The released pipeline has six stages: data download, pool construction with the attack generator, representation caching, the stream experiments, table generation and the latency benchmark. Every number in this paper is read from files that these stages write; a single shell script reruns them in order, and fixed seeds (pool seed 20261019, stream seeds 0--9) reproduce the pools and the stream orderings exactly. Training and inference are deterministic on CPU up to floating-point summation order. Code, per-run results and the manuscript source are available under the MIT licence at \url{https://github.com/osmancancet/slide-llm-phishing-drift}.
